\section{Research question and search scope}
\label{sec:review}
\textbf{Question.} Can a calibrated, structured uncertainty set around a process
reward model (PRM) reduce the gap between optimized proxy reward and genuine
reasoning quality? The comparison is with strong existing pessimistic methods
at matched compute, annotation budget, and policy movement.

We re-searched literature on process supervision, reward hacking, reward confidence
sets, pessimistic RLHF, distributional PRMs, selection-conditional conformal
prediction, and robust optimization. Primary sources are linked in the references.
Recent searches include September 2026 work and years distinguish initial preprints
from verified conference publication where available.

\subsection{What is already known, and what might be new}
Static PRM accuracy and resistance to optimization are different properties.
A policy or search procedure selects favorable prediction errors. Existing
defenses include better supervision, reward shaping, adversarial examples,
ensembles, confidence regions, and distributional rewards. Uncertainty around
rewards and pessimistic optimization are established ideas. The
potential contribution lies in the \emph{coverage target, geometry, and
optimization regime} for process rewards. 

Four close directions must shape our claims:
\begin{enumerate}
\item AdvPO already constructs reward confidence regions from last-layer features
and optimizes pessimistically \cite{advpo}.
\item PRM calibration and distributional PRMs already produce prefix-level
success estimates and uncertainty bounds \cite{park,otprm}.
\item Ensembles and reward-distribution DRO already give conservative effective
rewards \cite{coste,unifying,rrlhf}.
\item Selection-conditional conformal inference already addresses coverage after
data-dependent selection \cite{jomi}. Applying it to PRMs is an adaptation.
\end{enumerate}

\section{Process supervision: what is being estimated?}
\paper{Lightman et al., Let's Verify Step by Step (2023)}{
Human labels supervise intermediate mathematical reasoning; PRM800K makes this
supervision accessible. Process-supervised reward models outperform outcome
supervision in the paper's candidate-selection experiments, and active learning
improves annotation efficiency \cite{lightman}. This establishes the value of
intermediate checks, and not immunity to an optimizing policy. PRM800K provides
training data and a step-correctness target. A best-of-$N$ result should not be
interpreted as an RL robustness guarantee.}

\paper{Wang et al., Math-Shepherd (2023 preprint; ACL 2024)}{
Automatically generated process supervision supports reranking and stepwise PPO,
reducing dependence on human labels \cite{shepherd}. Completion-based labels
must be interpreted carefully as continuation success depends on the generator.
A confidence interval around that quantity need not cover intrinsic logical
correctness. The uncertainty target must be fixed before choosing data.}

\paper{Zhang et al., The Lessons of Developing PRMs (2025)}{
The paper contrasts Monte Carlo labels, LLM judgments, and human labels and
introduces consensus filtering. A correct answer can conceal flawed reasoning,
and a verifier optimized for best-of-$N$ can behave like an outcome verifier
\cite{lessons}. This motivates separate final-answer and process-error metrics.
Qwen PRMs are practical baselines, but a minimum score is not itself calibrated
uncertainty.}

\paper{Xie et al., Noise-aware Learning (2026); Ding et al., SCAN (2025)}{
Noise-aware learning uses reflection detection and iterative label refinement;
SCAN combines selective Monte Carlo annotation with denoising \cite{noise,scan}.
Both address corrupted synthetic supervision. Distinguish finite-rollout noise,
completer dependence, and reward-model uncertainty. Better labels may improve
the tube's center without characterizing coverage under optimization. Train
competing uncertainty methods on the same corrected labels.}

\paper{Pala et al., Error Typing for Smarter Rewards (2025)}{
PathFinder-PRM separates mathematical and consistency error prediction from
step-correctness scoring \cite{pathfinder}. Its reported 400K training samples
should not be called 400K complete trajectories without checking the data unit.
Error categories offer uncertainty features and stratified coverage diagnostics.
Richer supervision is distinct from an explicit set and from resistance to an
optimizing policy.}

\subsection{Very recent developments: error location and recovery}
\paper{Wang et al., Hierarchical Multi-Step Reward Models (2025; revised 2026)}{
HRM learns from individual and consecutive reasoning steps, while hierarchical
node compression augments tree-generated supervision \cite{hrm}. The method
motivates modeling multi-step coherence and recovery. This is not evidence that
ordinary prefix-conditioned PRMs ignore all history. A structured tube should
be compared with improving the reward representation itself, especially when
its claimed benefit concerns inter-step dependence.}

\paper{Han et al., Cliff (September 2026 preprint)}{
An off-the-shelf teacher locates the first mistake; Cliff assigns different
advantages before and after that boundary \cite{cliff}. This is a recent
credit-assignment comparator and suggests uncertainty over error locations.
Its effectiveness depends on teacher judgments. First-error localization does
not automatically resolve self-correction or teacher exploitation.}

\paper{Gan et al., Reasoning State Propagation (September 30, 2026 preprint)}{
RSP models transitions between valid and invalid reasoning states, including
break and repair probabilities, linking process labels to outcomes \cite{rsp}.
A weakest-step objective treats errors as irreversible; reasoning may recover.
An extension could put uncertainty around break/repair transitions, requiring
its own novelty comparison and more complex inference.}

\paper{Fan et al., TIPS (September 29, 2026 preprint)}{
Outcome-only RL trains a generative verifier to emit reasoning, step labels, and
an outcome label \cite{tips}. This is relevant to supervision costs and
generative-PRM baselines. Outcome rewards alone do not certify intermediate
labels, so audit them separately. These late preprints were screened at abstract
level; their gains are not independently reproduced evidence.}

\section{How process rewards are hacked}
\paper{Cheng et al., Stop Summation / PURE (NeurIPS 2025)}{
PURE uses minimum-based credit assignment to reduce accumulation of favorable
step errors \cite{pure}. It is a required baseline: an apparent uncertainty benefit
could actually result from aggregation. Reported mitigation is not a universal
guarantee against degenerate outputs. Test tubes with aggregation fixed, then
cross methods with sum, mean, minimum, and product.}

\paper{Zheng et al., CoLD (2025; revised May 2026)}{
CoLD targets length bias through a length adjustment, a learned bias estimator,
and joint training informed by counterfactual reasoning \cite{cold}. Length and
paraphrase transformations can probe our tube and motivate bias directions.
A fitted estimator is not an identified causal effect without assumptions.}

\paper{Tiwari et al., Reward Under Attack (2026)}{
Static perturbations, adversarial optimization, and RL expose PRM weaknesses.
In the studied AIME setup, the authors report near-perfect PRM scores while
answer accuracy stays below 4\% \cite{attack}. PRM-BiasBench and released code
are useful resources. This is evidence about specific models and settings, not
every PRM. Attack our robust score adaptively, not only the original reward.}

\paper{Juneja et al., Adversarial Training for PRMs (ICML 2026)}{
APRM co-trains an error generator and verifier to expose new failures
\cite{aprm}. It is a strong hard-negative comparator. Its oracle and game
formulation require assumptions: own-player strong concavity alone does not
imply strong monotonicity in general-sum games, and neural PPO parameters do
not automatically inherit policy-space concavity. We use the empirical method
without importing a blanket neural-training convergence guarantee.}

\paper{Pathmanathan and Huang, REFORM (2025 preprint; ACL 2026)}{
Reward-guided controlled decoding discovers falsely scored responses and uses
them to repair preference reward models \cite{reform}. This motivates acquiring
labels at optimizer-discovered errors. Compare tube construction against spending
the same labels on hard-negative retraining; otherwise improved data alone
could explain a gain. Its preference setting differs from process verification.}

\paper{Wallace et al., Universal Adversarial Triggers (2019)}{
Optimized token sequences expose learned NLP vulnerabilities \cite{triggers}.
They motivate suffix and discrete-token attacks, but do not establish transfer
to a particular PRM. Distinguish content attacks from parser or delimiter exploits.}

\section{Reward hacking theory and optimization pressure}
\paper{Skalse et al., Defining and Characterizing Reward Hacking (2022)}{
Hackability is defined through policy-order reversals between proxy and true
returns. Rich policy classes severely restrict nontrivial proxies that never
reverse preferences \cite{skalse}. This does not prove every training run hacks.
It motivates conditional guarantees over restricted or reachable policy classes
rather than globally perfect alignment.}

\paper{Gao et al., Scaling Laws for Reward Model Overoptimization (2022)}{
Experiments with a stronger gold reward model show proxy scores increasing after
gold scores deteriorate \cite{gao}. Use optimization-pressure curves, not one
checkpoint. The gold model remains a proxy for human intent; exact mathematics
and independent step audits are preferable where available. Empirical scaling
laws are not universal laws for PRM reasoning.}

\paper{Laidlaw et al., Correlated Proxies (2024 preprint; ICLR 2025)}{
Proxy--true correlation under a reference policy supports reward-hacking theory
and occupancy-measure regularization \cite{correlated}. Reference-distribution
correlation is distinct from pointwise validity in new states. For our project,
ask how a reward-set assumption yields a useful bound under changing reasoning
prefix distributions.}

\paper{Liu, Sun, and Zheng, Robust Optimization with Correlated Proxies (ICLR 2026)}{
The paper optimizes against worst-case rewards constrained by proxy correlation
and studies known linear reward features \cite{robustcorr}. Neither correlation
sets nor generic max--min hacking defenses are new. A PRM contribution needs
process-specific geometry, construction from supervision, or guarantees and
evidence under selection and on-policy shifts.}

\paper{Liu et al., Robust General Utility for RL (August 2026 preprint)}{
Utility functionals of occupancy are uncertain; algorithms address different
utility regimes \cite{generalutility}. This is relevant to nonlinear PRM
aggregation. It was screened at abstract level: inspect the full assumptions
before using its stationarity results. Generic robust utility optimization is
already prior work.}

\section{Closest prior work on uncertainty and pessimism}
\label{sec:closest}
\paper{Coste et al., Reward Model Ensembles (2023 preprint; revised 2024)}{
Mean, worst-case, and uncertainty-weighted ensembles are tested under best-of-$N$
and PPO, including noisy labels \cite{coste}. Pessimism reduces overoptimization
in the tested synthetic reward setting. Shared data or representations can make
members confidently wrong together. Include minimum-member and variance-penalty
baselines; simply using a lower ensemble score is not a new contribution.}

\paper{Zhang et al., AdvPO (2024)}{
AdvPO derives uncertainty from last-layer embeddings, constructs a confidence
region, and performs adversarial policy optimization \cite{advpo}. This is the
closest predecessor to a feature-based ellipsoidal tube. Theory assumes reward
approximation and feature structure; its practical radius is tuned. Potential
advances are process-level joint coverage, calibrated misspecification, and
selected-trajectory audits. Transferring its geometry to a PRM alone is a limited
extension.}

\paper{Yan et al., Reward-Robust RLHF (2024)}{
Bayesian reward ensembles and a nominal--pessimistic tradeoff already model
uncertain rewards \cite{rrlhf}. An ensemble distribution is a modeling device,
not automatic coverage of truth. Check coverage, sharpness, and systematic bias
independently.}

\paper{Park et al., Know What You Don't Know (2025)}{
Quantile regression calibrates PRM-derived continuation-success probabilities
for a specified generator. Bounds support adaptive inference budgets
\cite{park}. This is directly relevant, but its target differs from intrinsic
step correctness. Calibration and efficiency do not automatically give joint
coverage across selected candidates or an RL hacking guarantee. A faithful
implementation is a required baseline.}

\paper{Ma et al., Distributional PRMs via Conditional Optimal Transport (May 2026)}{
Conditional OT learns a monotone conditional quantile map using PRM features
and rollout-success targets \cite{otprm}. It offers coherent quantiles and flexible
confidence levels. Quantile monotonicity differs from finite-sample coverage
after selection; its target also changes with the completer. Our advance must
be more specific than learning reward distributions or confidence bounds.}

\paper{Yu et al., From Curiosity to Caution (ICLR 2026)}{
Caution penalizes prediction-error novelty to mitigate best-of-$N$ reward
hacking, with empirical and simplified theoretical analysis \cite{caution}.
Novelty is not logical error: unusual correct reasoning may be penalized while
familiar wrong reasoning evades it. Compare calibrated residual uncertainty
against novelty penalties at matched scoring cost.}

\paper{Hahami et al., A Unifying Lens on Reward Uncertainty (June 2026)}{
A distributional reward yields an entropic pessimistic score; mean, minimum,
and variance penalties arise as limits or approximations \cite{unifying}.
Reward robustness and policy KL have distinct roles. Entropic aggregation is
already prior work. Joint process-step distributions or policy-dependent
coverage could extend it, provided model misspecification is addressed.}

\paper{Zhou et al., PRISM (ACL 2026); Duan et al., BNRM (2026 preprint)}{
PRISM uses Gaussian mixtures to represent preference heterogeneity and
uncertainty; BNRM uses Bayesian nonnegative latent factors \cite{prism,bnrm}.
These concern preference rewards rather than certified process correctness.
Rich output distributions need not cover reasoning failures. Separate preference
diversity, annotation variability, and epistemic errors.}

\section{Statistical tools for constructing a meaningful set}
\paper{Lakshminarayanan et al., Deep Ensembles (2017)}{
Independently trained predictors provide practical uncertainty estimates
\cite{deepensemble}. Seeds, data resamples, and backbones give different diversity.
Measure systematic common errors. Dispersion is an input to calibration, not a
confidence-set theorem by itself.}

\paper{Abbasi-Yadkori et al., Linear Stochastic Bandits (2011)}{
Self-normalized martingale bounds support confidence ellipsoids for adaptive
features under a linear conditional mean and noise assumptions \cite{abbasi}.
Changing neural features, logistic labels, or corrupted annotations require
new arguments. A frozen-feature regression setting isolates these assumptions.}

\paper{Romano et al., Conformalized Quantile Regression (2019)}{
Conformal calibration corrects learned quantiles to give marginal prediction
coverage under exchangeability \cite{cqr}. It covers future observed targets,
not automatically conditional means or whole reward functions. Calibrate by
problem rather than treating correlated steps as independent examples.}

\paper{Jin and Ren, Confidence on the Focal (2024 preprint; JRSS B 2025)}{
Selection-conditional inference covers focal test units chosen by data-dependent
rules, including optimization-based selection \cite{jomi}. This is close to PRM
best-of-$N$. An adaptation must check rule symmetry, target, and computational
cost. It does not automatically cover an indefinitely co-adapting policy and
verifier.}

\paper{Tibshirani et al., Covariate Shift (2019); Barber et al., Beyond Exchangeability (2022)}{
Weighted and nonexchangeable conformal methods address specified shifts
\cite{weighted,beyond}. Density ratios, support overlap, and a stable conditional
label mechanism are substantive assumptions. Estimated or clipped weights do
not preserve exact coverage automatically. Changing the completer changes a
rollout-success target, so that is not mere covariate shift.}

\paper{Angelopoulos et al., Conformal Risk Control (2022)}{
Calibration extends to expected monotone losses \cite{crc}. This motivates
controlling false acceptance of bad steps. The theorem must match the loss and
monotonicity: selecting a different trajectory at each threshold can break
monotonicity. Start with thresholded accept/reject rules on fixed scores.}

\paper{Bao et al., Model Selection for Conformalized Robust Optimization (2025)}{
Prediction sets, model selection, and downstream robust decision quality are
connected \cite{croms}. A narrow valid interval is not necessarily best for a
decision. Decision-aware calibration is prior work; report accuracy and regret
alongside width.}

\subsection{Late statistical results to monitor}
Two 2026 preprints were screened: a testing view of conformal universality and
impossibility \cite{conformaltest}, and finite-sample CQR bounds including
fixed-score calibration under known shift \cite{cqrlimits}. They may inform
coverage--efficiency lower bounds. Full assumptions need inspection; no theorem
below is attributed to these abstracts.

\section{Robust optimization and reward shaping}
\paper{Duchi and Namkoong, Uniform Performance via DRO (2018 preprint)}{
Distributional ambiguity emphasizes difficult examples and tail behavior, with
statistical tradeoffs \cite{duchi}. Robust PRM training can reweight ordinary and
discovered failures, but empirical reweighting cannot create unseen reasoning
behaviors. A set of data distributions differs from a set of reward functions.}

\paper{Iyengar, Robust Dynamic Programming (2005)}{
Rectangular transition uncertainty supports robust dynamic programming
\cite{iyengar}. The structural lesson transfers: independent local adversarial
choices and one globally consistent model are different problems. Text
transitions may be known while PRM reward uncertainty is coupled. Do not assume
a robust Bellman recursion for an arbitrary nonrectangular tube.}

\paper{Ng, Harada, and Russell, Policy Invariance (1999)}{
Potential-based shaping preserves policy preferences under appropriate boundary
conditions \cite{shaping}. Where outcome verification is reliable, a PRM can
assist learning as a potential. This differs from replacing rewards with
pessimistic estimates. It does not remove verifier errors or promise faster
learning.}

\section{Evaluation resources and evidence gaps}
ProcessBench tests earliest-error identification on expert-annotated mathematical
solutions \cite{processbench}; PRMBench tests dimensions including simplicity,
soundness, and sensitivity \cite{prmbench}; Socratic-PRMBench adds reasoning-pattern
diagnostics \cite{socratic}. PRM-BiasBench supplies controlled perturbations
\cite{attack}. Static performance cannot replace optimization stress tests.

\begin{longtable}{p{0.20\linewidth}p{0.35\linewidth}p{0.36\linewidth}}
\toprule
Family & Established capability & What our study must resolve \\
\midrule
\endhead
Process supervision & Intermediate verification & What quantity the tube covers \\
Label denoising & Better synthetic labels & Coverage after optimization \\
PURE / shaping & Better credit assignment or outcome invariance & Tube benefits
with aggregation held fixed \\
Ensembles / AdvPO & Pessimistic reward optimization & Shared bias and joint
process uncertainty \\
Quantile / OT PRMs & Prefix distributions and bounds & Selection coverage and
completer dependence \\
Conformal / JOMI & Guarantees for specified sampling rules & Feasible PRM
adaptation and adaptive RL assumptions \\
Correlated-proxy robustness & Structured worst-case rewards & Construction from
process supervision \\
\bottomrule
\end{longtable}

\subsection{A defensible candidate contribution}
\emph{Optimization-aware uncertainty for process rewards}: compare pointwise
intervals with structured or problem-level tubes; measure coverage after
selection; determine whether coverage improvements reduce hacking at matched
optimization pressure. Novelty remains provisional. An empirical paper could
expose systematic post-selection undercoverage and a practical repair. A theory
paper could characterize safe improvement and failure of marginal intervals.
A combined paper would connect a restricted theorem to measured violations in
real PRMs. No publication outcome is assumed.

\subsection{Reading order}
Read Part II, then Lightman and the Qwen lessons for targets; PURE and Reward
Under Attack for failures; AdvPO, Park, conditional OT, Caution, and reward
uncertainty unification for close methods; JOMI and correlated-proxy robustness
for statistical distinctions. Part III develops concrete designs.
